% !TEX root = ../main.tex
\section{Métodos}
\label{sec:methods}

La \cref{fig:ejemplo} muestra una figura de ejemplo y la \cref{tab:ejemplo} una tabla.
Las figuras van en la carpeta figures/ con nombres en minúscula y sin espacios.

\begin{figure}[htbp]
  \centering
  \includegraphics[width=0.6\linewidth]{example-image-a} % replace with figures/<name>
  \caption{Figura de ejemplo.}
  \label{fig:ejemplo}
\end{figure}

\begin{table}[htbp]
  \centering
  \caption{Tabla de ejemplo con unidades.}
  \label{tab:ejemplo}
  \begin{tabular}{l S[table-format=1.3]}
    \toprule
    Magnitud & {Valor (\si{\electronvolt})} \\
    \midrule
    $E_0$ & 0.125 \\
    $E_1$ & 0.375 \\
    \bottomrule
  \end{tabular}
\end{table}
